\chapter{Distribuciones muestrales}
\labch{distribuciones-muestrales}

\section{La ley de los grandes números}

En la clase hicimos un pequeño experimento tirando una moneda diez veces. No recuerdo los resultados, pero aquí lo hice unas cuantas veces
\begin{table}[!h]
\centering
\begin{tabular}[t]{rrS}
\toprule
Número de Caras & Frecuencia & Proporción\\
\midrule
1 & 1 & 0.05\\
2 & 1 & 0.05\\
3 & 1 & 0.05\\
4 & 6 & 0.30\\
5 & 7 & 0.35\\
6 & 3 & 0.15\\
7 & 1 & 0.05\\
\bottomrule
\end{tabular}
\end{table}

La mayoría de los resultados están dentro del intervalo \([4, 6]\), lo que es más o menos esperable. Si nos metemos en el intervalo \([3, 7]\) ya están casi todos los resultados: sólo faltan dos.

Sin embargo, si hacíamos el experimento 1000, 5000 o más veces la intuición nos decía que esa diferencia tenía que irse achicando, ¿no? NO.

Nuestra intuición no está tan errada, pero hay que afinarla. En general, cuando tenemos una suma (o promedio o conteo) tendremos que \[\text{suma} = \text{valor esperado} + \text{error aleatorio.}\] La pregunta es qué tan grande es ese error aleatorio. Hay una fórmula que nos permite establecer el tamaño de este error, llamado \emph{error estándar.} Cuando contabilicemos la suma de las extracciones de una caja, con reposición, el error estándar de la suma es \[\sqrt{\text{número de extracciones}} \cdot (\text{desviación estándar de la caja})\]

En otras palabras, el error \emph{aumenta} a medida que repetimos nuestro experimento, de manera proporcional a la raíz del número de repeticiones. En ese sentido, nuestra intuición estaba completamente incorrecta. Sin embargo, si consideramos el promedio \[\frac{\sqrt{\text{número de extracciones}} \cdot (\text{SD de la caja})}{\text{número de extracciones}} = \frac{\text{SD de la caja}}{\sqrt{\text{número de extracciones}}} \]
y cuando el número de extracciones crece, el error estándar decrece. 

Esta idea está representada en la \emph{ley de los grandes números.} Hay varias versiones de esta ley, pero todas son teoremas que hablan de la convergencia de una serie de variables aleatorias.

\formal
Si \(X_1, X_2, \cdots, X_n\) son variables aleatorias independientes e idénticamente distribuidas, con esperanza \(\mu\) y varianza \(\sigma^2\) finitas, entonces la variable \(\mean{X} = \sum_{i=1}^{n} X_i/n\) cumple que para todo \(\varepsilon >0\), \[\lim_{n\to\infty} \prob\left(|\mean{X} - \mu| \leq \varepsilon\right) = 0\]